\chapter{Competitor Rules}

\section{Equipment}
\subsection{Unicycles}
Only regular unicycles may be used.
The maximum outer diameter of the wheel is 640\,mm (24+ Class).
In addition, the unicycles must not have sharp or protruding parts anywhere that might cause injuries.
This refers especially to quick-release levers and bolts.
The pedals must be plastic or rubber.

\subsection{Rider Identification}
All players of a team must wear shirts of the same color.
The color must be clearly different from the opponent's color.
At tournaments and other large events each team should have two different colored sets of shirts.
Players must display identifying numbers.

\subsection{Sticks}
All sticks legal for playing ice-hockey or floorball (apart from those for the goalkeeper) can be used.
Cracked or splintered sticks must be taped or repaired before play.
An upper end made of rubber is recommended.

\subsection{Safety}
Attention must be drawn to the safety of the players and spectators.
Thus, the rules must be obeyed strictly and all equipment must be in good condition.

All items that protrude from the body that may cause injury (for example watches, necklaces, earrings) must be removed.
In instances where this is impossible, the items must be covered sufficiently to remove likelihood of injury.
Shoes must be worn and shoelaces must be short or tucked in.
The following optional clothing is suggested: knee pads, gloves, helmets, safety glasses and dental protection.

\section{Players and Team Officials}
A team consists of at least three players and may include one registered Head Coach.
A team on the field consists of up to five players with a team requiring a minimum of three players to begin a match.
Each player can be the goalkeeper at any time.
The goalkeeper has no special rights.

\subsection{Captain}
One Captain shall be appointed by each team, and they and the Head Coach alone shall be authorized to consult with the referees regarding any queries on rule interpretation.
Change of the team captain shall only take place in case of injury, illness or penalty box for the remainder of the game.
The team captain shall wear an armband, which shall be worn on the upper arm and be clearly visible.

\subsection{Player Registration}
A player may only play for one team.
If a tournament includes multiple grades of competition (e.g. A, B, C), the player may only play within a single grade.

\subsection{Player Substitutions}
Player substitutions are possible at any time.
During active play, the substituting players must interchange at the same location and within their own half.
The player entering the field may only enter the field after the substituting player has completely left it.
It is not necessary to indicate substitutions to the referee.

\section{Event Flow}
\subsection{Game Duration}
The play time is given by the playing schedule and is a relative play time.
The time stops only at the request of the referee.
The teams change sides during the break.
At the start of each period, all players must be in their own half of the field.
Each period starts with a face-off at the center mark.
If the game ends in a draw and a decision is necessary, play is continued with extended time.
If it's still a draw, a decision is reached with a penalty shootout.

\subsection{Penalty Shootout}\label{subsec:hockey_penalty_shootout}
Three of the players from each team get one penalty shot each.
If it is still a draw, each team shoots one more penalty until there is a decision.
It is possible that one player can take more than one shot.
However, in all cases at least two other players have to take a shot before the same player can shoot again.

For the penalty, all players except for the defending goalkeeper leave the corresponding half of the playing field.
The referee places the ball on the center point and indicates that play may start with a whistle.
The goalkeeper must remain close to the goal line until the attacking player has made contact with the ball.
After the referee's whistle, the attacking player may play the ball, trying to score a goal.
The shot attempt begins as soon as the player taking the shot has contacted the ball, and the player must remain in motion towards the goal line with no backwards movement or stopping allowed.

Once the ball has been shot, the play shall be considered complete.
No goal can be scored on a rebound of any kind (an exception being the ball off the goal post and/or the goalkeeper and then directly into the goal), and any time the ball crosses the goal line, the shot shall be considered complete.

\subsection{Riding The Unicycle}
The player has to be riding the unicycle freely.
A player may only use the stick as support.
Touching or holding the wall or goal to gain an advantage over your opposition is illegal. 
A player who is falling off the unicycle may take part in the game until touching the ground.
A remounting player must sit on the seat and have both feet on the pedals before participating in the game again.
If a player who is not riding a unicycle shoots into their own goal, the advantage rule applies for the attacking team and the goal is valid.

\subsection{Contact With The Ball}
The stick, the unicycle and the whole body can be used to play the ball.
It all counts as a contact.
Players are allowed to play the ball with the body twice in a row only if one of the contacts is passive.
When the ball is played with the body, the player must not catch or otherwise hold the ball and the contact with the ball should be instantaneous.
For arms and hands see also section \ref{subsec:hockey_goal-shots_with-arms-or-hands}.

\subsection{Start and Stop}
Starting and resuming the game is always initiated by the referee's whistle.
If a team starts to play before the referee's whistle, it is stopped immediately by two or more quick consecutive blows of the whistle.
Then, the previous referee ruling is repeated.
When the referee blows the whistle during the game, it is interrupted immediately.

\subsection{Restart After A Goal}
After a goal, the non-scoring team gets the ball.
All players must go to their own half.
After the referee's whistle, the game resumes when the ball or a player of the team in possession crosses the center line.
It is legal to directly shoot a goal after passing the center line, for example without passing the ball to another player first.

\subsection{Ball Out Of Bounds}
If the ball leaves the field, or touches an object above the field, the game is interrupted immediately, regardless of whether the ball returns to play.
The team opposite to that of the player who last touched the ball gets a free shot (for execution see \ref{subsec:hockey_free_shot} Free Shot).

\subsection{Ball In Spokes}
In the event that a ball becomes lodged within a player's wheel, possession will be given to the opposing team as a free shot (For execution see \ref{subsec:hockey_free_shot} Free Shot).

\subsection{Time-Outs}
During each knockout game (e.g. half finals, finals), each team shall have the right to request one 60-second time-out.

The team captain must request the time-out by informing the first referee during a stoppage in play.
Upon receiving the request, the first referee will signal the time-out to the scoring desk.
The secretary will record which team has taken the time-out, and the timekeeper will measure 60\,seconds before play resumes.
After a time-out, play shall be resumed according to what caused the interruption.

A penalised player is not allowed to participate in a time-out.

\section{Goal Shots}
\subsection{Conditions for Awarding a Goal}
A goal is considered scored when the entire ball crosses the rear edge of the goal line, provided it entered from the front of the goal and was played legally with no offence committed by the attacking team immediately before or during the goal.

If the goal is displaced, a goal shall be awarded if the ball enters the displaced goal and the displacement was caused by the defending team, or if the ball crosses the marked goal line between the original positions of the goalposts and beneath the imaginary crossbar.

\subsection{Goal Shot With Arms Or Hands \label{subsec:hockey_goal-shots_with-arms-or-hands}}
A goal is disallowed if scored with arms or hands.
The defending team gets a free shot (goalkeeper's ball).
This rule does not apply if the ball is shot into one's own goal.

\subsection{Long Shot}
For a goal to be valid, the last contact with the ball must occur in the opponent’s half, meaning the entire ball has crossed the center line into the opponent’s half.
Otherwise it's considered a Long Shot and the  defending team gets a free shot (goalkeeper's ball).
This rule does not apply if the ball is shot from the opponents' half into one's own goal.

\subsection{Ball In The Outside Of The Net}
If the ball becomes lodged in the outside of the goal net, or if the ball entered the goal through a hole in the back or side of the net, a free shot is given against the team whose player last played the ball.

\section{Fouls}
\subsection{General Considerations}
All players must take care not to endanger others.
Exaggerated roughness can lead to injuries and must therefore be avoided.
The game is non-contact: the opponents and their unicycles may not be touched.

\subsection{Right Of Way}
The following rules apply when riders come into contact with each other:
\begin{itemize}
\item No player may endanger another player by forcing them to give way (for example, to push them toward the wall).
\item A player who is idling or resting on the stick must be evaded.
  However, the idling or resting player must ensure the stick does not SUB players as per rule \ref{subsec:hockey_fouls_sub}.
\item The leading of two players riding next to each other may choose the direction of turns.
If both are evenly side-by-side, the one in possession of the ball may choose the direction.
\item If two players are approaching each other directly or at an obtuse angle, both must take care to avoid contact.
  If contact occurs, the referee will penalise the player deemed to have caused the contact.
\item In all cases not mentioned above, it is up to the referee to make a decision.
\end{itemize}

\subsection{SUB (Stick Under Bike) \label{subsec:hockey_fouls_sub}}
A player who holds their stick in a way that someone else rides over or against it is always committing a foul regardless of the situation.

\subsection{SIB (Stick In Bike)}
If a stick gets into the spokes of an opponent, the holder of the stick is committing a foul.

\subsection{Insults}
A player must not insult the referee or other players.

\subsection{Moving The Goal}
Players must not move the goal out of position.
If the goal is displaced, the offending team must return it to its correct position.

If the non-offending team has possession, play continues under advantage (as the infringement is ongoing).
If the goal is replaced before possession changes, play continues without interruption.
However, if possession changes before the goal is restored, play is stopped and the non-offending team is awarded a free shot (corner or goalkeeper’s ball).

If the offending team has possession when the goal is moved, play is stopped and the non-offending team is awarded a free shot (corner or goalkeeper’s ball).

If it is unclear which team moved the goal, play is stopped and resumed with a face-off at the nearest corner mark.

\subsection{Obstacle}
A player who is off the unicycle must not be an obstacle for opponents.
The player is considered an obstacle if the player, the unicycle or stick is hit by the ball and also if an opponent cannot move around freely.
The player should remount at the same spot, but if necessary move out of the way of play first.

\subsection{Throwing Sticks}
A player must not intentionally drop or throw their stick.

\subsection{Top Of The Stick}
The upper end of the stick must always be covered with one hand to avoid injury to other players.
A brief removal of the upper hand from the stick to play the ball with that hand is acceptable provided that this is done in a safe manner.

\subsection{High Stick}
The entire blade of the stick must always be below the players' own hips and the hips of all players in the vicinity who might be endangered.
Exception: When defending a shot on goal in the direct vicinity of one's own goal, the lower end of the stick can be raised as high as the crossbar of the goal.

\subsection{Stick Contact}
Only in the vicinity of the ball (defined as the ball within the radius of the outstretched arm length plus stick) may a player touch an opponent's stick with their stick to block them.
However, this contact may not be hard.
The players must take care not to hit an opponent with their stick, especially after a shot.

\subsubsection{Hooking}
It is illegal to use the blade of the stick in a manner that enables the player to restrain an opponent's stick.

\subsubsection{Slashing}
Any downward or forceful chop with the stick on an opponent’s stick that, in the judgment of the referee, involves excessive force, shall be penalized as slashing.
Raising an opponent’s stick is permitted unless it involves excessive force, such as causing the opponent’s stick to reach high-stick height.
In such cases, it shall be considered slashing.

\subsection{Intentional Delay Of Game}\label{subsec:hockey_delay_game}
Intentional delay of the game shall result in a penalty and may also lead to a stoppage of time.
A delay of game foul shall be called when a team intentionally delays the restart of play.
Examples of delay of game include deliberately hitting the ball out of bounds, taking excessive time to execute a free shot, slowly returning to their own half after scoring a goal, or deliberately delaying the restart after conceding a goal. 

\section{Penalties}
In instances of a violation of the rules, the referee must penalize the offending team or play the advantage.
However, rule violations that do not influence the course of the game shall not be penalized.
Any rule violation resulting in a player being sent to the penalty box, (2-minute, 5-minute, or match send-offs) shall be deemed to have influenced the course of the game.
When playing the advantage the referee does not blow the whistle but should display the hand sign for a free shot and shout ``Advantage!''
In the event that an advantage was not gained, the referee should enforce the appropriate penalty from the initial point of infringement or, when the penalty has occurred within the goal area, the closest corner mark or 2\,m in front of goal line.
In cases where an infringement occurs outside the goal area and within 2\,m of the sideline, the free shot shall be given from a spot located 2\,m from the sideline.
Additionally, if advantage was gained before a send-off was issued (see \ref{subsec:hockey_penalties_player_send-off} Player Send-off), the referee shall resume play with a face-off at the location of the ball when play was stopped.
When two or more players fall and/or it is unclear whether a foul occurred, the referees can interrupt the game and restart it with a face-off.

\subsection{Free Shot}\label{subsec:hockey_free_shot}
The free shot is the standard penalty for all violations of the rules, and where a ball leaves the playing field.
It is applied in all cases except for those explicitly mentioned in sections \ref{subsec:hockey_penalties_65m}-\ref{subsec:hockey_penalties_face-off}.
For penalty violations, the free shot is executed from the point where the violation was done.
If the ball is ruled out of bounds, the free shot is taken from the point where the ball was judged to have left the field, 2\,m in from the sidelines, or below the point that it contacted an object above the field.
Exceptions: If a team receives a free shot within the opponents' goal area, the free shot is executed at the closest corner mark (corner shot).
If a team receives a free shot within their own goal area, the free shot is taken at a distance of 2\,m in front of the goal line (goalkeeper's ball).

If a delay of game penalty is awarded, the continuation of the game depends on the previous situation.
\begin{itemize}
	\item If the delay of game occurs during the execution of a free shot, the opposing team will receive a free shot.
	\item If the delay of game is caused by the scoring team, play will resume as normal with a restart after the goal.
	\item If the delay of game is caused by the non-scoring team, play will continue with a free shot from the centre mark.
\end{itemize}

The free shot is indirect.
The player executing the free shot may only touch the ball once until a contact by another player occurs.
The ball shall be hit with the stick, not dragged, flicked or lifted on the stick.
Opposing players must keep a distance with their unicycles and their sticks of at least 2\,m from the ball.

\subsection{6.5\,M \label{subsec:hockey_penalties_65m}}
If legal play would have led to a direct chance to score a goal, a ``6.5\,m'' is given.
The following situations are a prevention of a direct chance to score and should be punished with a 6.5\,m penalty:
\begin{itemize}
\item An attacking player is fouled in the opposition goal area while in a strong position to score.
\item An attacking player is fouled when moving towards the opposition goal with a single opponent in front.
\end{itemize}
The ball is placed at the 6.5\,m mark.
A player of the defending team goes to the goal and must sit with the bottom of the wheel of their unicycle within 1\,m of the goal line.

The attacking team shall designate one player to take the 6.5\,m shot.
All other players must exit the goal area.
Upon the referee’s whistle, all players except the shooter must ensure that both their unicycles and sticks remain at least 2\,m away from the 6.5\,m mark until the ball has been struck.

The 6.5\,m is direct.
The shooter may only touch the ball once.
The ball shall be hit with the stick, not dragged, flicked or lifted on the stick.
If no goal is scored, play continues as soon as the ball touches a goal post, the goal's crossbar, the goalkeeper touches the ball or the ball crosses the extended goal line.
At this moment, the shooting player may play the ball again.
A 6.5\,m awarded at the end of, or after a time period has ended, is still executed but play does not continue after an unsuccessful shot.

\subsection{Penalty Goal}\label{subsec:hockey_penalties_penalty_goal}
If the defending team prevents a goal from being scored through an illegal play and if, in the opinion of the referee, the ball was traveling directly toward the goal and would definitely have entered the goal without being touched by another player, a penalty goal may be awarded to the attacking team.
If there is any doubt as to the certainty of a goal, a 6.5\,m must be awarded as described in section \ref{subsec:hockey_penalties_65m}.

\subsection{Face-off \label{subsec:hockey_penalties_face-off}}
To resume the game without penalizing one of the teams, a face-off can be used.

A face-off during the game is executed where the ball was when the game was interrupted but at least 2\,m from the side line.
Exception: Within the goal area, the face-off is executed at the closest corner mark.

One player from each team may take part in the face-off with all other players' unicycles and sticks at a distance of at least 2\,m from the ball.
In a face-off, each player must begin with the blade of their stick on the floor on the side of the face-off point nearest to their ground line.
This side is determined by a line through the face-off point that runs parallel to the center line.
For the face-off, the referee drops the ball between two opposing players.
The ball should be dropped from 50\,cm (approx. 20\,inch).
Play starts when the ball touches the ground as signalled by the executing referees whistle.

\subsection{Player Send-off}\label{subsec:hockey_penalties_player_send-off}
The referee can send a player off the field for two minutes, five minutes, or for the remainder of the game.

If a player from the team in possession of the ball commits an offence warranting a 2-min, 5-min, or match penalty, the referee shall immediately interrupt the game and impose the penalty.
If a player from the team not in possession commits the offence, play shall continue until the offending team gains control of the ball, at which point play is stopped (including stoppage of timer) and the penalty assessed.

If the penalty to be imposed is a 2-min penalty and a goal is scored during the continued play, the 2-min penalty shall not be imposed.
This does not apply to 5-minute or match penalties, which shall be imposed regardless of whether a goal is scored.
If two or more 2-min penalties are to be imposed and a goal is scored during the continued play, the captain of the offending team shall designate which penalty(ies) are to be assessed and which one is cancelled due to the goal.

In instances of a penalty, the offending team must play with one fewer than their legal maximum for the duration of the penalty, or until the end of the current match, whichever comes first (see \ref{subsubsec:hockey_penalties_multiple_penalties} Multiple penalties).
Any unexpired penalty at the end of regular time shall remain in effect during extra time and also any penalty shootout (\ref{subsec:hockey_penalty_shootout}) if applicable.
All penalties except match penalties shall be considered terminated only after the final result of the match is determined, including any penalty shootout if applicable.

If the referees are unable to identify the offender, (e.g. incorrect substitution), the team captain shall choose a non-penalised player to serve the penalty.
If the team captain cannot provide a player or is themselves penalised, the referees shall choose the player.

Penalized players are prohibited from communicating with their team for the duration of the penalty.
A penalised team captain may not communicate with the referees unless directly addressed by them.
When enforcing a penalty, the referee should signal to stop the timer to discuss the required punishment, provide a precise start time for the penalty, and to explain their ruling to players.

\subsubsection{2-Minute Penalty}
While a player is in the penalty box for a 2-minute penalty, the team may not substitute a replacement for that player.

The following list of offences lead to a 2-minute penalty; however, this list is not exhaustive.
Referees may impose a 2-minute penalty for other actions that, in their judgment, constitute unsporting behaviour, create a safety risk, or significantly violate the spirit of fair play.
\begin{itemize}
\item When a player intentionally commits a foul.
\item When a player commits a SUB or SIB at speed.
\item When a player forces or pushes an opponent into the walls or the goal causing them to fall.
\item When a player is guilty of dangerous play or careless physical play.
\item When a player hits another player with their stick above hip height.
\item When a player aggressively swings their stick at an opponent or the opponents stick (“slashing”), regardless of contact.
Non-aggressive stick contact below knee height should not lead to a 2-minute penalty, but may still lead to a free shot.
Raising an opponent’s stick above hip height is a foul, but does not automatically warrant a 2-minute penalty.
\item When a player throws their stick.
\item When a player violates right-of-way rules while riding at speed.
If the referee deems that the infraction is not solely attributed to one player, or it was not clearly avoidable, only a free shot shall be given.
\item When a player persistently backtalks to the referee or questioning their decisions.
\item When a player insults the referees, players or spectators.
\item When a player actively obstructs the execution of a free shot.
This includes violations of the 2\,m rule during a free shot, corner shot, goalkeeper's ball, 6.5\,m, or the restart after a goal (i.e. crossing the centre line before an opposing player or the ball has crossed the centre line).
\item When a player uses incorrect equipment or a player is missing correct rider identification (e.g. no identifying number).
Where possible, the referee should inform players of equipment and identification issues prior to the start of the match.
\item When an incorrect substitution takes place.
If the substitution occurs in a borderline or near-simultaneous manner, a penalty shall only be given if it impacts the play.
\item When a team has too many players on the field.
The last player to enter the field shall receive a 2-minute penalty.
The offending team must then reduce the number of players on the field so that they are playing with one fewer than their permitted number.
\item When a team systematically disrupts play by committing repeated offences leading to a free shot.
This also includes when a team commits a number of rule violations during a short time.
The player committing the last rule violation shall serve the penalty.
\item When a team intentionally delays play.
The referee shall issue a warning to the captain on the first offence.
A 2-minute penalty shall be applied for any further intentional delay.
\end{itemize}

\subsubsection{5-Minute Penalty}
While a player is in the penalty box for a 5-min penalty, the team may not substitute a replacement for that player.

The following list of offences lead to a 5-minute penalty; however, this list is not exhaustive.
\begin{itemize}
\item Repeated fouls by a player who has previously received a 2 minute penalty
\item Intentional dangerous foul
\item Violent conduct against other players, their team officials or spectators
\end{itemize}

\subsubsection{Match Penalty}
When a player is sent off for the remainder of the game, they may not take part in the current match and their team's following match.
However, in the current match, the penalised team may bring a player on after a five minute period.
%TODO The following sentences could be interpreted to mean that the penalized player is removed from the game and, in addition, another player must leave the field and go to the penalty box for 5 minutes—but I believe the intention is that a substitute player goes to the penalty box
The offending player is removed from the match but does not enter the penalty box.
A non-penalised player, selected by the offending team’s captain, must serve the five-minute penalty in the penalty box.
If the captain is unable to select a player, the referees shall choose one.
Once the five-minute penalty has expired, that player may return to the field.

The following list of offences lead to a match penalty; however, this list is not exhaustive.
\begin{itemize}
\item Repeated fouls by a player who has previously received a 5 minute penalty
\item Repeated violence of a player who has already received 5 minutes before
\item Violence against referees
\end{itemize}

\subsubsection{Multiple penalties}\label{subsubsec:hockey_penalties_multiple_penalties}
A team may receive any number of 2-minute, 5-minute or match penalties during a game, and individual players may be given multiple 2-minute or 5-minute penalties, including while already serving a previous penalty.
No player may serve more than one penalty simultaneously.
If a player who is already serving a penalty in the penalty box receives an additional penalty, the penalties shall be served consecutively, not simultaneously.
The timer for the first penalty continues without interruption, and the second penalty begins only once the first has expired.
If the enforcement of a penalty would result in the team having no remaining players on the field, the team shall forfeit the match.

\subsubsection{Termination}
A team is considered ``short-handed'' when it has fewer players on the field than its opponent due to one or more 2-min penalties.
If a team is ``short-handed'' and the opposing team scores a goal or is awarded one (\ref{subsec:hockey_penalties_penalty_goal} Penalty Goal), the earliest active 2-min penalty automatically terminates.
If a team already short-handed incurs an additional 2-min penalty and the non-offending side scores during continued play, the goal counts, the new penalty is assessed, and the earliest existing 2-min penalty terminates.
If a 2-minute penalty is to be imposed against a team already serving a 5-minute or match penalty, and the non-offending team scores during continued play, the 2-minute penalty shall not be assessed due to the scoring of the goal.

\section{Protests}
Protests must be filed on an official form within two hours of the posting of event results.
Every effort will be made for all protests to be handled within 30 minutes from the time they are received.
